\chapter{Method}
\label{cha:method}
In this chapter, a step by step description of the process of detecting anomalies in a raw unstructured log file is provided. 
\section{Framework}
To perform anomaly detection on log data, the data has to be transformed and prepared to fit the selected anomaly detection techniques. The used method of choice consists of four steps: 

\begin{enumerate}
    \item \textbf{Log collection}: Has already been completed and is provided in the selected datasets, this study focuses on the remaining steps.
    \item \textbf{Log parsing}: Structures the log files.
    \item \textbf{Feature extraction}: Creates numerical vectors based on the structured data.
    \item \textbf{Anomaly detection}: Uses the vectors created as input for the algorithms to detect anomalies.
\end{enumerate}
An overview of this process is shown in Figure \ref{fig:method-process}.

For the offline log data evaluations, the entire log file will be handled at once. That is, the entire log file will first be parsed and structured, followed by feature extraction and finally, anomaly detection will be performed on the entire dataset at once. In contrast, for the online log data evaluation, the entire log file will still be parsed at once, as parsing each row individually does not impact the results. However, feature extraction will be performed dynamically, processing one structured log row at a time. Each extracted feature vector will be immediately analyzed for anomalies before moving on to the next log row. 


\begin{figure}[t]
    \centering
    \includegraphics[width=1\linewidth]{figures/Method.pdf}
    \caption[Methodology overview of our framework]{Methodology overview of our framework. Created in BioRender \cite{biorendermethodology}.}
    \label{fig:method-process}
\end{figure}


\section{Log parsing}
Logs are generated by various systems and often contain valuable information about the system. However, these logs are stored in raw text formats, making it difficult to process. The huge volume logs combined with the unstructured representation makes it impractical to manually inspect \cite{zhu2019toolsbenchmarksautomatedlog}.  Log parsing is therefore essential to prepare the data in a structured way for the anomaly detection algorithms, making it the first step in log data analysis. By extracting the key attributes, going from free text raw-log messages to structured form, log parsing enables the possibility to use automated anomaly detection techniques to detect patterns on log data.

The log messages are generated by a logging statement in the source code, and consists of a message header and message content. The logging framework has determined the header and is rather easy to extract (e.g. of headers are timestamps and level). It is however harder to structure the free-text message content, because it's a formation of constant strings and variable values \cite{zhu2019toolsbenchmarksautomatedlog}. The goal of log parsing is to transform these messages into structured event templates by identifying the core event type while filtering out dynamic variables. This abstraction enables more efficient log analysis and anomaly detection.

A presentation, analyses and evaluation is performed by Zhu et al \cite{zhu2019toolsbenchmarksautomatedlog} on $13$ different log parsers on $16$ datasets. They go through each algorithm, evaluating based on accuracy, robustness and efficiency. The top performing algorithm was shown to be Drain, providing the highest accuracy on average, but shows smallest variance combined with robustness when changing volumes of logs. Pinja et al \cite{Pinja2017Drain} also performed an analysis, evaluating Drain on five real world log data sets. The results show that Drain demonstrate the highest accuracy on four of the sets, and a comparable accuracy on the fifth while also improving running time by $51.8$\textasciitilde $81.47$\%.  Additionally, most log parsing methods focus only on offline batch processing. However, Drain is a method suitable for both online and offline log parsing \cite{Pinja2017Drain}. Therefore, Drain is the method selected for the HDFS data. Due to time limitations, we used an already parsed BGL dataset \cite{4273008} \cite{zhu2023loghublargecollectionlog}.

\subsection{Drain}
Drain is a method supporting both online and offline log parsing that only requires the raw log messages \cite{Pinja2017Drain}. Drain can automatically extract log messages, split it into log groups using a parse tree with fixed depth, (i.e. all leaf nodes are found at the same depth). This is done by preprocessing the log message by regular expressions. Then searching for a log group (leaf nodes) by following the rules of the internal nodes of the tree. When a suitable group is found, the message is matched with that log event, else a new log group is created. By using a parse tree with fixed depth, the technique avoids comparing log messages to every log group, saving time and efficiency. The log group contains a log event and \revrev{an}{a} log ID, the first describes the log message of the group (the message content described above). The second, log ID, records the ID of log messages in that group. In this project, an existing implementation by the LogPai team was used.


\section{Feature extraction}
Once the parsed log messages are created, the structured log messages need to be converted into numerical features for machine learning algorithms to be able to process them. Therefore, a feature extraction is needed to encode the parsed log messages further. This can be done using grouping techniques, such as fixed window, sliding windows and session windows. By doing so, feature vectors are created representing the number of occurrences for a specific event. All feature vectors can together form a feature matrix \cite{Shilin2017experiencereport}. 

\subsection{Fixed Window}
\label{fixedWindowSection}
Fixed window is a timestamped based method. The logs are divided into windows based on the fixed time duration, where the window size $\Delta t$ is a constant value \cite{Shilin2017experiencereport}. Therefore, the number of windows depend on the predefined window size, and the windows do not overlap with each other's time boundaries \cite{Gautam2018bigdatarealtime}. Further, logs that occur in the same windows are called a log sequence. A demonstration of fixed window is provided in Figure \ref{fig:Fixed-window}.

\begin{figure}[t]
    \centering
    \includegraphics[width=1\linewidth]{figures/Fixed_window.pdf}
    \caption[Fixed window approach]{Fixed window approach, where the \textcolor{orange}{orange} and \textcolor{blue}{blue} circles represent different types of event. Created in BioRender \cite{biorenderfixedwindow}.}
    \label{fig:Fixed-window}
\end{figure}



\subsection{Sliding Window}
Unlike fixed windows, sliding windows are defined by two key parameters, window size and step size, where the windows overlap if the window size is greater than step size \cite{Gautam2018bigdatarealtime}. In Figure \ref{fig:sliding-window} it is shown that the logs that occur in the same sliding window are grouped as a log sequence. Logs may also be duplicated in several sliding windows due to the overlap that occurs \cite{Shilin2017experiencereport}. 

\begin{figure}[t]
    \centering
    \includegraphics[width=1\linewidth]{figures/Slidingwindow.pdf}
    \caption[Sliding window approach]{Sliding window approach, where the \textcolor{orange}{orange} and \textcolor{blue}{blue} circles represent different types of event. Created in BioRender \cite{biorenderslidingwindow}.}
    \label{fig:sliding-window}
\end{figure}

\subsection{Session Window}
The session window approach is however based on identifiers rather than timestamps. These identifiers help to differentiate execution paths in log datasets \cite{Shilin2017experiencereport}. For example, the HDFS logs use \textit{block\_id} to track operations as allocation, writing, replication or deletion of blocks. By leveraging these identifiers the logs can be grouped into session windows, where each session window is associated with an identifier. This is demonstrated in Figure \ref{fig:session-win}, where the number of event types are grouped \revrev{in}{by} sessions. 

\begin{figure}[t]
    \centering
    \includegraphics[width=1\linewidth]{figures/Session window.pdf}
    \caption[Session window approach]{Session window approach, where the \textcolor{orange}{orange} and \textcolor{blue}{blue} circles represent different types of event occurring in one session. Created in BioRender \cite{biorendersessionwindow}.}
    \label{fig:session-win}
\end{figure}

\section{Anomaly detection}
Once the vector matrix is constructed, the process of detecting anomalies can begin. This study employs machine learning-based approaches, specifically leveraging \textbf{PCA}, \textbf{LogClustering}, \textbf{Isolation Forest}, \textbf{Half-Space Trees} and \textbf{One-Class SVM} as the selected algorithms. A detailed description of these methods is provided in Section \ref{AD}. These are all unsupervised machine learning-based methods (described in Section \ref{ML}), as they align most closely with real-world scenarios where labeled data is unavailable.  \revrev{}{Additionally, \textbf{SVM} was implemented using supervised training, this result serves as a benchmark for the unsupervised methods. This allows an estimation of an upper bound of achievable performance when labels are 
available, and compare how well unsupervised methods perform. }

\revrev{}{The methods were selected to represent different categories of anomaly detection approaches, such as tree-based, clustering based, and statistical techniques. This was done to ensure a broad perspective in our evaluation. Additionally, the chosen algorithms have demonstrated promising results in previous research, which further motivated their inclusion in this study.} 
For PCA, Isolation Forest and LogClustering, open source implementations by the LogPai team were used \cite{Shilin2017experiencereport}.  We implemented the Half-Space Trees and One-Class SVM algorithms using the River library in Python \cite{montiel2021river}, where a stochastic implementation of the One-Class SVM algorithm has been created to handle online scenarios.



\section{Implementation}
In our experiments, we used $80$\% of the dataset as a training set and remaining $20$\% as a test set for the online data. This provides enough data for both training and testing, while correlating with a real world scenario. \revrev{}{ Additionally, the online methods trained on normal data exclusively.} For the offline data however, the entire file was used as training, this also correlates with a real world situation where one would have access to the entire file. The HDFS dataset was first log parsed using Drain, then feature extracted with session windows based on \textit{block\_id} before implementing the different anomaly detectors. The online implementations on the HDFS dataset were initially trained in an offline setting, as access to the entire training data reflects a real-world scenario. The testing was then implemented in an online environment. For the BGL dataset, because of time limitations an already log parsed file in a structured format was used \cite{zhu2023loghublargecollectionlog} \cite{4273008}\revrev{}{, since log parsing with Drain would provide the same results}. Since there is no feature in the BGL dataset that can be used to create session windows, sliding windows was used as a feature extraction method, since it performs better on BGL than fixed window \cite{Shilin2017experiencereport}. The offline sliding window size was timestamp-based while the online sliding window size was based on the number of log-rows, to simplify computation. Both BGL implementation scenarios were implemented with an event count vector. Lastly, for One-Class SVM a standard scaler was used after the feature extraction to preprocess the data for better results \cite{montiel2021river}. For Half-Space Trees, a MinMaxScaler was applied for the same purpose. \revrev{}{The LogClustering algorithm was trained only on normal data, unlike the other two offline algorithms which used both normal and anomalous data in the training phase. While SVM trained in a supervised setting, with a labeled dataset to serve as a benchmark.} The code was written in Python 3.11.6 on a macOS system running on a MacBook Air with an Apple M2 chip ($8$-core CPU, $8$-core GPU) and $16$ GB unified memory. 